\documentclass{article}
\usepackage{amsmath}
\usepackage{amssymb}
\usepackage{geometry}
\usepackage{graphicx}
\usepackage{fancyhdr}
\usepackage{enumitem}
\usepackage{xcolor}

\geometry{a4paper, margin=1in}

\title{\textbf{Homework 14.4: Tangent Planes and Linear Approximations}}
\author{Tashfeen Omran}
\date{February 2026}

\begin{document}

\maketitle
\tableofcontents
\newpage

\section{Part 1: Comprehensive Introduction}

\subsection{Core Concepts}
This topic extends the fundamental concept of the \textbf{tangent line} from single-variable calculus to multivariable calculus. In Calculus I, a differentiable function $y = f(x)$ can be approximated locally by a tangent line. In Calculus III, for a function of two variables $z = f(x, y)$, the graph is a surface in 3D space. If the surface is smooth (differentiable) at a point, we can approximate it locally with a flat plane known as the \textbf{Tangent Plane}.

Associated with the tangent plane is the concept of \textbf{Linearization}. The equation of the tangent plane provides a linear function $L(x, y)$ that approximates the actual complex function $f(x, y)$ for points near the point of tangency. This is crucial for estimation and error analysis.

We also introduce the \textbf{Total Differential}, $dz$, which represents the approximate change in the function value ($z$) resulting from small changes in the input variables ($dx$ and $dy$).

\subsection{Intuition and Derivation}
\textbf{The Tangent Plane:}
Imagine zooming in infinitely close to a point on a smooth curved surface (like a sphere or a hill). As you zoom in, the surface begins to look flat. This "flatness" is the tangent plane.

To derive the equation, recall that a plane is determined by a point $P(x_0, y_0, z_0)$ and a normal vector $\mathbf{n}$, or alternatively, by two vectors lying in the plane.
1. Consider the curves formed by intersecting the surface with the planes $y = y_0$ and $x = x_0$.
2. The slope of the tangent line to the first curve is the partial derivative $f_x(x_0, y_0)$. A vector in this direction is $\mathbf{v_1} = \langle 1, 0, f_x(x_0, y_0) \rangle$.
3. The slope of the tangent line to the second curve is $f_y(x_0, y_0)$. A vector in this direction is $\mathbf{v_2} = \langle 0, 1, f_y(x_0, y_0) \rangle$.
4. The normal vector $\mathbf{n}$ is the cross product $\mathbf{v_1} \times \mathbf{v_2} = \langle -f_x, -f_y, 1 \rangle$.
5. Using the dot product formula for a plane, we arrive at the standard equation.

\subsection{Historical Context}
The development of tangent planes and differentials was part of the broader 18th and 19th-century effort to rigorousize calculus for physical applications. Mathematicians like Lagrange and Cauchy sought to understand how complex systems (like fluid flow or gravitational fields) behave locally. The problem they solved was: "How can we replace a complicated, non-linear system with a simple linear one to make calculations possible?" This is the foundation of modern engineering simulations, where curved objects are modeled as meshes of tiny flat triangles (tangent planes).

\subsection{Key Formulas}
\begin{enumerate}
    \item \textbf{Equation of the Tangent Plane} at $(x_0, y_0, z_0)$:
    \[ z - z_0 = f_x(x_0, y_0)(x - x_0) + f_y(x_0, y_0)(y - y_0) \]
    
    \item \textbf{Linearization} $L(x, y)$ at $(a, b)$:
    \[ L(x, y) = f(a, b) + f_x(a, b)(x - a) + f_y(a, b)(y - b) \]
    
    \item \textbf{Total Differential} $dz$:
    \[ dz = \frac{\partial z}{\partial x} dx + \frac{\partial z}{\partial y} dy = f_x(x, y) dx + f_y(x, y) dy \]
    
    \item \textbf{Linear Approximation for 3 Variables} $w = f(x, y, z)$:
    \[ L(x, y, z) = f(a,b,c) + f_x(a,b,c)(x-a) + f_y(a,b,c)(y-b) + f_z(a,b,c)(z-c) \]
\end{enumerate}

\subsection{Prerequisites}
To master this topic, you must be proficient in:
\begin{itemize}
    \item \textbf{Partial Derivatives:} calculating $f_x$ and $f_y$ using power, product, quotient, and chain rules.
    \item \textbf{Equation of a Plane:} Understanding point-normal form (though we derive a specific version here).
    \item \textbf{Vector Algebra:} Dot products and basic vector components.
    \item \textbf{Algebraic Evaluation:}Accurately plugging numbers into complex expressions.
\end{itemize}

\newpage
\section{Part 2: Detailed Homework Solutions}

\subsection*{Problem 1: Tangent Plane to a Paraboloid}
\textbf{Question:} Find an equation of the tangent plane to the surface $z = 23 - x^2 - y^2$ at the specified point. (Note: The point is not explicitly listed in the text, but based on the grid and standard problem variations, we will use the point $(2, 3)$. If your specific numbers differ, substitute them into the general form derived below.)

\textbf{Solution:}
Let the function be $f(x, y) = 23 - x^2 - y^2$.
We assume the point of tangency is where $x=2$ and $y=3$.

\textbf{Step 1: Find the z-coordinate of the point.}
\[ z_0 = f(2, 3) = 23 - (2)^2 - (3)^2 = 23 - 4 - 9 = 10 \]
The point is $P(2, 3, 10)$.

\textbf{Step 2: Calculate Partial Derivatives.}
\[ f_x(x, y) = \frac{\partial}{\partial x}(23 - x^2 - y^2) = -2x \]
\[ f_y(x, y) = \frac{\partial}{\partial y}(23 - x^2 - y^2) = -2y \]

\textbf{Step 3: Evaluate Partial Derivatives at the point $(2, 3)$.}
\[ f_x(2, 3) = -2(2) = -4 \]
\[ f_y(2, 3) = -2(3) = -6 \]

\textbf{Step 4: Apply the Tangent Plane Formula.}
\[ z - z_0 = f_x(x_0, y_0)(x - x_0) + f_y(x_0, y_0)(y - y_0) \]
\[ z - 10 = -4(x - 2) - 6(y - 3) \]

\textbf{Step 5: Simplify.}
\[ z - 10 = -4x + 8 - 6y + 18 \]
\[ z = -4x - 6y + 26 + 10 \]
\[ z = -4x - 6y + 36 \]
Or in standard linear form: $4x + 6y + z = 36$.

\textbf{Final Answer:} \[ z = 36 - 4x - 6y \]

---

\subsection*{Problem 2: Tangent Plane Setup (Step-by-Step)}
\textbf{Question:} Let $z = f(x, y) = 4x^2 + y^2 - 7y$. Find $f_x(x, y)$, $f_y(x, y)$, evaluate at $(1, 3, -8)$, and find the tangent plane equation.

\textbf{Solution:}
\textbf{Step 1: Find Partial Derivatives.}
\[ f_x(x, y) = \frac{\partial}{\partial x}(4x^2 + y^2 - 7y) = 8x \]
\[ f_y(x, y) = \frac{\partial}{\partial y}(4x^2 + y^2 - 7y) = 2y - 7 \]

\textbf{Step 2: Evaluate at the point $(1, 3)$.}
Here $x=1, y=3$.
\[ f_x(1, 3) = 8(1) = 8 \]
\[ f_y(1, 3) = 2(3) - 7 = 6 - 7 = -1 \]

\textbf{Step 3: Tangent Plane Equation.}
Point: $(1, 3, -8)$. Slopes: $8$ and $-1$.
\[ z - (-8) = 8(x - 1) + (-1)(y - 3) \]
\[ z + 8 = 8x - 8 - y + 3 \]
\[ z = 8x - y - 5 - 8 \]
\[ z = 8x - y - 13 \]

\textbf{Final Answer:}
$f_x(x,y) = 8x$ \\
$f_y(x,y) = 2y - 7$ \\
Equation: \[ z = 8x - y - 13 \]

---

\subsection*{Problem 3: Tangent Plane Variation}
\textbf{Question:} Find the tangent plane to $z = 2x^2 + y^2 - 7y$ at $(1, 3, -10)$.

\textbf{Solution:}
\textbf{Step 1: Partial Derivatives.}
\[ f_x = 4x \]
\[ f_y = 2y - 7 \]

\textbf{Step 2: Evaluate at $(1, 3)$.}
\[ f_x(1, 3) = 4(1) = 4 \]
\[ f_y(1, 3) = 2(3) - 7 = -1 \]

\textbf{Step 3: Equation.}
\[ z - (-10) = 4(x - 1) - 1(y - 3) \]
\[ z + 10 = 4x - 4 - y + 3 \]
\[ z = 4x - y - 1 - 10 \]
\[ z = 4x - y - 11 \]

\textbf{Final Answer:} \[ z = 4x - y - 11 \]

---

\subsection*{Problem 4: Exponential Function}
\textbf{Question:} Find the tangent plane to $z = e^{x-y} + 8$ at $(1, 1, 9)$.

\textbf{Solution:}
\textbf{Step 1: Partial Derivatives (Chain Rule).}
\[ f_x = \frac{\partial}{\partial x}(e^{x-y} + 8) = e^{x-y} \cdot (1) = e^{x-y} \]
\[ f_y = \frac{\partial}{\partial y}(e^{x-y} + 8) = e^{x-y} \cdot (-1) = -e^{x-y} \]

\textbf{Step 2: Evaluate at $(1, 1)$.}
\[ f_x(1, 1) = e^{1-1} = e^0 = 1 \]
\[ f_y(1, 1) = -e^{1-1} = -e^0 = -1 \]

\textbf{Step 3: Equation.}
\[ z - 9 = 1(x - 1) - 1(y - 1) \]
\[ z - 9 = x - 1 - y + 1 \]
\[ z - 9 = x - y \]
\[ z = x - y + 9 \]

\textbf{Final Answer:} \[ z = x - y + 9 \]

---

\subsection*{Problem 5: Product Rule with Exponential}
\textbf{Question:} Find the tangent plane to $z = y^2 e^x$ at $(0, 6, 36)$.

\textbf{Solution:}
\textbf{Step 1: Partial Derivatives.}
\[ f_x = y^2 \frac{\partial}{\partial x}(e^x) = y^2 e^x \]
\[ f_y = e^x \frac{\partial}{\partial y}(y^2) = 2y e^x \]

\textbf{Step 2: Evaluate at $(0, 6)$.}
\[ f_x(0, 6) = (6)^2 e^0 = 36(1) = 36 \]
\[ f_y(0, 6) = 2(6) e^0 = 12(1) = 12 \]

\textbf{Step 3: Equation.}
\[ z - 36 = 36(x - 0) + 12(y - 6) \]
\[ z - 36 = 36x + 12y - 72 \]
\[ z = 36x + 12y - 36 \]

\textbf{Final Answer:} \[ z = 36x + 12y - 36 \]

---

\subsection*{Problem 6: Trigonometric Product Rule}
\textbf{Question:} Find the tangent plane to $z = x \sin(y - x)$ at $(7, 7, 0)$.

\textbf{Solution:}
\textbf{Step 1: Partial Derivatives.}
For $f_x$, use Product Rule: $u = x$, $v = \sin(y-x)$.
\[ f_x = (1) \cdot \sin(y-x) + x \cdot \cos(y-x) \cdot \frac{\partial}{\partial x}(y-x) \]
\[ f_x = \sin(y-x) + x \cos(y-x)(-1) \]
\[ f_x = \sin(y-x) - x \cos(y-x) \]

For $f_y$, $x$ is constant:
\[ f_y = x \cdot \cos(y-x) \cdot \frac{\partial}{\partial y}(y-x) \]
\[ f_y = x \cos(y-x)(1) = x \cos(y-x) \]

\textbf{Step 2: Evaluate at $(7, 7)$.}
Note that the argument inside trig functions is $y - x = 7 - 7 = 0$.
\[ f_x(7, 7) = \sin(0) - 7 \cos(0) = 0 - 7(1) = -7 \]
\[ f_y(7, 7) = 7 \cos(0) = 7(1) = 7 \]

\textbf{Step 3: Equation.}
\[ z - 0 = -7(x - 7) + 7(y - 7) \]
\[ z = -7x + 49 + 7y - 49 \]
\[ z = -7x + 7y \]

\textbf{Final Answer:} \[ z = -7x + 7y \]

---

\subsection*{Problem 7: Logarithmic Function}
\textbf{Question:} Find the tangent plane to $z = \ln(x - 3y)$ at $(4, 1, 0)$.

\textbf{Solution:}
\textbf{Step 1: Partial Derivatives.}
\[ f_x = \frac{1}{x - 3y} \cdot \frac{\partial}{\partial x}(x - 3y) = \frac{1}{x - 3y} \]
\[ f_y = \frac{1}{x - 3y} \cdot \frac{\partial}{\partial y}(x - 3y) = \frac{-3}{x - 3y} \]

\textbf{Step 2: Evaluate at $(4, 1)$.}
Denominator: $x - 3y = 4 - 3(1) = 1$.
\[ f_x(4, 1) = \frac{1}{1} = 1 \]
\[ f_y(4, 1) = \frac{-3}{1} = -3 \]

\textbf{Step 3: Equation.}
\[ z - 0 = 1(x - 4) - 3(y - 1) \]
\[ z = x - 4 - 3y + 3 \]
\[ z = x - 3y - 1 \]

\textbf{Final Answer:} \[ z = x - 3y - 1 \]

---

\subsection*{Problem 8: Linear Approximation (3 Variables)}
\textbf{Question:} Find the linear approximation of $f(x, y, z) = \sqrt{x^2 + y^2 + z^2}$ at $(2, 9, 6)$ and approximate $\sqrt{2.03^2 + 8.99^2 + 5.98^2}$.

\textbf{Solution:}
\textbf{Step 1: Value at the center point $(2, 9, 6)$.}
\[ f(2, 9, 6) = \sqrt{2^2 + 9^2 + 6^2} = \sqrt{4 + 81 + 36} = \sqrt{121} = 11 \]

\textbf{Step 2: Partial Derivatives.}
\[ f_x = \frac{1}{2\sqrt{x^2+y^2+z^2}} \cdot 2x = \frac{x}{\sqrt{x^2+y^2+z^2}} \]
Similarly by symmetry:
\[ f_y = \frac{y}{\sqrt{x^2+y^2+z^2}}, \quad f_z = \frac{z}{\sqrt{x^2+y^2+z^2}} \]

\textbf{Step 3: Evaluate Partials at $(2, 9, 6)$.}
Denominator is 11.
\[ f_x = \frac{2}{11}, \quad f_y = \frac{9}{11}, \quad f_z = \frac{6}{11} \]

\textbf{Step 4: Linear Approximation Formula.}
\[ L(x, y, z) = 11 + \frac{2}{11}(x - 2) + \frac{9}{11}(y - 9) + \frac{6}{11}(z - 6) \]

\textbf{Step 5: Approximate the specific value.}
Target: $x = 2.03, y = 8.99, z = 5.98$.
Calculate differentials ($\Delta$ values):
\[ \Delta x = 2.03 - 2 = 0.03 \]
\[ \Delta y = 8.99 - 9 = -0.01 \]
\[ \Delta z = 5.98 - 6 = -0.02 \]

Substitute into $L$:
\[ L \approx 11 + \frac{2}{11}(0.03) + \frac{9}{11}(-0.01) + \frac{6}{11}(-0.02) \]
\[ L \approx 11 + \frac{1}{11} [ 2(0.03) + 9(-0.01) + 6(-0.02) ] \]
\[ L \approx 11 + \frac{1}{11} [ 0.06 - 0.09 - 0.12 ] \]
\[ L \approx 11 + \frac{1}{11} [ -0.15 ] \]
\[ L \approx 11 - \frac{0.15}{11} \approx 11 - 0.013636... \]
\[ L \approx 10.9864 \]

\textbf{Final Answer:} \[ 10.9864 \]

---

\subsection*{Problem 9: Step-by-Step 3-Variable Approximation}
\textbf{Question:} Same function $f(x, y, z) = \sqrt{x^2 + y^2 + z^2}$, center $(9, 6, 2)$. Approximate $\sqrt{9.01^2 + 5.98^2 + 1.97^2}$.

\textbf{Solution:}
\textbf{Step 1: Partials.}
As derived above:
\[ f_x(x, y, z) = \frac{x}{\sqrt{x^2+y^2+z^2}}, \quad f_y = \frac{y}{\sqrt{\dots}}, \quad f_z = \frac{z}{\sqrt{\dots}} \]

\textbf{Step 2: Evaluate at $(9, 6, 2)$.}
\[ \text{Value} = \sqrt{81 + 36 + 4} = \sqrt{121} = 11 \]
\[ f_x = \frac{9}{11}, \quad f_y = \frac{6}{11}, \quad f_z = \frac{2}{11} \]

\textbf{Step 3: Linear Approximation Form.}
\[ f(x, y, z) \approx 11 + \frac{9}{11}(x - 9) + \frac{6}{11}(y - 6) + \frac{2}{11}(z - 2) \]

\textbf{Step 4: Approximation.}
Target: $(9.01, 5.98, 1.97)$.
\[ \Delta x = 0.01, \quad \Delta y = -0.02, \quad \Delta z = -0.03 \]
\[ \Delta f \approx \frac{9(0.01) + 6(-0.02) + 2(-0.03)}{11} \]
\[ \Delta f \approx \frac{0.09 - 0.12 - 0.06}{11} = \frac{-0.09}{11} \]
\[ \text{Result} = 11 - \frac{0.09}{11} = 11 - 0.008181... \approx 10.9918 \]

\textbf{Final Answer:} \[ 10.9918 \]

---

\subsection*{Problem 10: Differential of a Log Function}
\textbf{Question:} Find the differential $dz$ for $z = x \ln(y^3 + 9)$.

\textbf{Solution:}
Formula: $dz = \frac{\partial z}{\partial x} dx + \frac{\partial z}{\partial y} dy$.
\[ \frac{\partial z}{\partial x} = \ln(y^3 + 9) \cdot 1 = \ln(y^3 + 9) \]
\[ \frac{\partial z}{\partial y} = x \cdot \frac{1}{y^3 + 9} \cdot \frac{\partial}{\partial y}(y^3 + 9) = x \frac{3y^2}{y^3 + 9} \]

\textbf{Final Answer:}
\[ dz = \left( \ln(y^3 + 9) \right) dx + \left( \frac{3xy^2}{y^3 + 9} \right) dy \]

---

\subsection*{Problem 11: Differential with Time Variable}
\textbf{Question:} Find the differential $dz$ for $z = e^{-9x} \cos(4\pi t)$.

\textbf{Solution:}
Formula: $dz = \frac{\partial z}{\partial x} dx + \frac{\partial z}{\partial t} dt$.
\[ \frac{\partial z}{\partial x} = (-9)e^{-9x} \cos(4\pi t) \]
\[ \frac{\partial z}{\partial t} = e^{-9x} (-\sin(4\pi t)) (4\pi) = -4\pi e^{-9x} \sin(4\pi t) \]

\textbf{Final Answer:}
\[ dz = \left( -9e^{-9x} \cos(4\pi t) \right) dx + \left( -4\pi e^{-9x} \sin(4\pi t) \right) dt \]

---

\subsection*{Problem 12: Differential of 3 Variables}
\textbf{Question:} Find $dH$ for $H = x^8 y^6 + y^4 z^7$.

\textbf{Solution:}
Formula: $dH = \frac{\partial H}{\partial x} dx + \frac{\partial H}{\partial y} dy + \frac{\partial H}{\partial z} dz$.
\[ \frac{\partial H}{\partial x} = 8x^7 y^6 \]
\[ \frac{\partial H}{\partial y} = x^8 (6y^5) + (4y^3) z^7 = 6x^8 y^5 + 4y^3 z^7 \]
\[ \frac{\partial H}{\partial z} = y^4 (7z^6) = 7y^4 z^6 \]

\textbf{Final Answer:}
\[ dH = (8x^7 y^6) dx + (6x^8 y^5 + 4y^3 z^7) dy + (7y^4 z^6) dz \]

---

\subsection*{Problem 13: Tangent Plane Fill-in-the-Blanks}
\textbf{Question:} Tangent plane to $z = 3x^2 + 5y^2$ at $(1, 1, 8)$.

\textbf{Solution:}
1. $f_x(x, y) = 6x$.
2. $f_y(x, y) = 10y$.
3. $f_x(1, 1) = 6(1) = 6$.
4. $f_y(1, 1) = 10(1) = 10$.
5. Tangent Plane: $z - 8 = 6(x - 1) + 10(y - 1)$.
6. Solve for $z$:
\[ z = 6x - 6 + 10y - 10 + 8 \]
\[ z = 6x + 10y - 8 \]

\textbf{Final Answer:}
Blanks: $6x$, $10y$, $6$, $10$.
Equation: $z - 8 = 6(x - 1) + 10(y - 1)$ or $z = 6x + 10y - 8$.

---

\subsection*{Problem 14: Linearization and Comparison}
\textbf{Question:} $f(x, y) = 5xe^{xy}$. Point $(1, 0)$. Estimate $f(1.1, -0.1)$.

\textbf{Solution:}
\textbf{Step 1: Partials.}
\[ f_x = 5(1 \cdot e^{xy} + x \cdot ye^{xy}) = 5e^{xy}(1 + xy) \]
\[ f_y = 5x(xe^{xy}) = 5x^2 e^{xy} \]

\textbf{Step 2: Evaluate at $(1, 0)$.}
\[ f(1, 0) = 5(1)e^0 = 5 \]
\[ f_x(1, 0) = 5e^0(1 + 0) = 5 \]
\[ f_y(1, 0) = 5(1)^2 e^0 = 5 \]

\textbf{Step 3: Linearization.}
\[ L(x, y) = 5 + 5(x - 1) + 5(y - 0) = 5 + 5(x - 1) + 5y \]

\textbf{Step 4: Approximate.}
$x = 1.1, y = -0.1$.
\[ L(1.1, -0.1) = 5 + 5(0.1) + 5(-0.1) = 5 + 0.5 - 0.5 = 5 \]

\textbf{Step 5: Compare with actual.}
Actual: $f(1.1, -0.1) = 5(1.1)e^{(1.1)(-0.1)} = 5.5 e^{-0.11}$.
$e^{-0.11} \approx 0.89583$.
$5.5(0.89583) \approx 4.92709$.

\textbf{Final Answer:}
$L(x,y) = 5 + 5(x-1) + 5y$ \\
Approximation: $5$ \\
Actual: $\approx 4.92709$

---

\subsection*{Problem 15: Differentiability and Linearization}
\textbf{Question:} $f(x, y) = x^3 y^2$ at $(-3, 2)$.

\textbf{Solution:}
\textbf{Step 1: Partials.}
\[ f_x = 3x^2 y^2 \]
\[ f_y = 2x^3 y \]

\textbf{Step 2: Evaluate at $(-3, 2)$.}
\[ f_x(-3, 2) = 3(-3)^2 (2)^2 = 3(9)(4) = 108 \]
\[ f_y(-3, 2) = 2(-3)^3 (2) = 2(-27)(2) = -108 \]
Function value: $f(-3, 2) = (-3)^3 (2)^2 = -27(4) = -108$.

\textbf{Step 3: Linearization.}
\[ L(x, y) = -108 + 108(x - (-3)) - 108(y - 2) \]
\[ L(x, y) = -108 + 108(x + 3) - 108(y - 2) \]
Simplify:
\[ L(x, y) = -108 + 108x + 324 - 108y + 216 \]
\[ L(x, y) = 108x - 108y + 432 \]

\textbf{Final Answer:}
$f_x = 3x^2 y^2, \quad f_y = 2x^3 y$ \\
$f_x(-3, 2) = 108, \quad f_y(-3, 2) = -108$ \\
$L(x, y) = 108x - 108y + 432$

\newpage
\section{Part 3: In-Depth Analysis}

\subsection{A) Problem Types and General Approach}
\begin{itemize}
    \item \textbf{Type 1: Tangent Plane Calculation (Problems 1-7, 13)}
    \textit{Strategy:} Calculate $z_0$, find partial derivatives $f_x$ and $f_y$, evaluate them at $(x_0, y_0)$, and plug everything into the point-normal plane equation.
    
    \item \textbf{Type 2: Linear Approximation \& Estimation (Problems 8, 9, 14, 15)}
    \textit{Strategy:} This is identical to Type 1, but we view the tangent plane equation as a function $L(x,y)$. The key additional step is calculating the "offsets" $dx = x - a$ and $dy = y - b$ and computing $dz = f_x dx + f_y dy$.
    
    \item \textbf{Type 3: Total Differential Computation (Problems 10-12)}
    \textit{Strategy:} These problems test the symbolic application of the chain rule. You do not evaluate at a point; you simply write the sum of partials multiplied by their respective differentials ($dx, dy, dt, dz$).
\end{itemize}

\subsection{B) Key Manipulations}
\begin{itemize}
    \item \textbf{Product Rule within Partials (Problem 6):} In $z = x \sin(y-x)$, finding $f_x$ required the product rule because $x$ appears both as a coefficient and inside the sine function.
    \item \textbf{Chain Rule for Logarithms (Problem 10):} For $z = x \ln(y^3 + 9)$, finding $f_y$ required the chain rule $\frac{d}{dy}(\ln u) = \frac{u'}{u}$.
    \item \textbf{Handling 3 Variables (Problem 8, 9, 12):} The formulas naturally extend by adding a term. $dz$ becomes $dw = f_x dx + f_y dy + f_z dz$.
    \item \textbf{Implicit Point Identification (Problem 1):} Sometimes the point is given on a grid. Recognizing integer intersections on a graph is a useful skill when text is missing.
\end{itemize}

\section{Part 4: "Cheatsheet" and Tips for Success}

\subsection{Formulas}
\begin{itemize}
    \item \textbf{Tangent Plane:} $z = z_0 + f_x(x-x_0) + f_y(y-y_0)$
    \item \textbf{Linearization:} $L(x,y) \approx f(a,b) + f_x(a,b)\Delta x + f_y(a,b)\Delta y$
    \item \textbf{Differential:} $dz = f_x dx + f_y dy$
\end{itemize}

\subsection{Tricks and Shortcuts}
\begin{itemize}
    \item \textbf{The "Delta" Trick:} When approximating values like $\sqrt{4.01^2 + 2.98^2}$, do not plug 4.01 into the plane equation. Instead, calculate just the \textit{change} ($\Delta z$) using $\Delta x = 0.01$ and $\Delta y = -0.02$, then add it to the "nice" value (5). It's much faster arithmetic.
    \item \textbf{Symmetry:} In symmetric functions like $\sqrt{x^2+y^2+z^2}$, once you find $f_x$, $f_y$ and $f_z$ will have the exact same form, just swapping variables.
\end{itemize}

\subsection{Pitfalls}
\begin{itemize}
    \item \textbf{Sign Errors in Subtraction:} Remember $\Delta x = \text{Target} - \text{Start}$. If you go from 3 to 2.9, $\Delta x = -0.1$, not $0.1$.
    \item \textbf{Forgetting the Chain Rule:} In $e^{x-y}$, $f_y = e^{x-y} \cdot (-1)$. That negative sign is often missed.
    \item \textbf{Mixing Variables:} In $dz = (\dots)dx + (\dots)dy$, ensure the $dx$ term contains the $x$-derivative and the $dy$ term contains the $y$-derivative.
\end{itemize}

\section{Part 5: Conceptual Synthesis}
\subsection{Thematic Connections}
The core theme is \textbf{Local Linearity}. Just as a curve looks like a line when you zoom in (Calculus I), a surface looks like a plane when you zoom in (Calculus III). This connects to the concept of \textbf{Manifolds} in advanced geometry—spaces that look Euclidean (flat) on small scales.

\subsection{Links}
\textbf{Backward:} Relies heavily on the equation of lines and planes (vectors) and partial derivatives.
\textbf{Forward:} This is the foundation for \textbf{Optimization}. Max/min points occur where the tangent plane is horizontal ($f_x=0, f_y=0$). It also leads to the \textbf{Total Derivative} matrix (Jacobian) in vector calculus.

\section{Part 6: Real-World Application}

\subsection{A) Concrete Scenarios}
\begin{enumerate}
    \item \textbf{Quantitative Finance (Option Pricing):} The Black-Scholes equation for pricing options relies on partial derivatives. The "Delta" of an option is exactly $f_S$ (partial derivative with respect to stock price), and "Theta" is $f_t$ (partial with respect to time). Traders use linear approximations to hedge risks (Delta Hedging).
    \item \textbf{Thermodynamics:} The Ideal Gas Law $PV = nRT$. If Pressure and Volume change slightly, engineers use differentials $dT = \frac{\partial T}{\partial P}dP + \frac{\partial T}{\partial V}dV$ to estimate the temperature change without recalculating the whole state.
\end{enumerate}

\subsection{B) Model Problem Setup (Finance)}
\textbf{Scenario:} A call option price $C$ depends on the stock price $S$ and time $t$.
\textbf{Model:} Let $C(S, t)$ be the price.
We know the current price $S_0$ and time $t_0$.
Equation to solve for change in price $dC$:
\[ dC \approx \frac{\partial C}{\partial S} dS + \frac{\partial C}{\partial t} dt \]
Traders use this to estimate how much money they lose if the stock drops by \$1 ($dS = -1$) or if one day passes ($dt = 1$).

\section{Part 7: Common Variations}
\textbf{Missing Concept: The Normal Line.}
The homework asked for tangent planes, but often exams ask for the \textbf{Normal Line}—the line perpendicular to the surface at a point.
\textbf{Method:} The normal vector to the plane is $\mathbf{n} = \langle f_x, f_y, -1 \rangle$.
\textbf{Example:} For $z = x^2 + y^2$ at $(1, 2, 5)$.
$f_x = 2, f_y = 4$. Normal vector $\mathbf{n} = \langle 2, 4, -1 \rangle$.
Parametric Equations:
\[ x = 1 + 2t, \quad y = 2 + 4t, \quad z = 5 - t \]

\section{Part 8: Advanced Diagnostic Testing}
\textit{Find the error in the following solutions.}

\textbf{Problem 1 (Flawed):} Find tangent plane to $z = x^2 y$ at $(1, 2, 2)$.
\textit{Flawed Solution:} $f_x = 2xy, f_y = x^2$. At $(1,2)$: $f_x = 4, f_y = 1$.
Equation: $z = 4(x-1) + 1(y-2)$.
\textbf{Error:} The $z_0$ term is missing on the LHS.
\textbf{Correction:} $z - 2 = 4(x-1) + 1(y-2)$.

\textbf{Problem 2 (Flawed):} Approximate $f(3.1, 4.2)$ for $f(x,y)=x y$ using $(3,4)$.
\textit{Flawed Solution:} $f_x = y = 4$, $f_y = x = 3$.
$L = 12 + 4(3.1) + 3(4.2)$.
\textbf{Error:} Multiplied slopes by the full values $x, y$ instead of the differences $\Delta x, \Delta y$.
\textbf{Correction:} $L = 12 + 4(0.1) + 3(0.2)$.

\textbf{Problem 3 (Flawed):} Differential of $z = e^{2x}$.
\textit{Flawed Solution:} $dz = 2e^{2x}$.
\textbf{Error:} Missing the $dx$ differential term.
\textbf{Correction:} $dz = 2e^{2x} dx$.

\textbf{Problem 4 (Flawed):} Tangent plane to $z = \sin(x) + y$.
\textit{Flawed Solution:} $f_x = \cos(x), f_y = 0$.
\textbf{Error:} Derivative of $y$ with respect to $y$ is 1, not 0.
\textbf{Correction:} $f_y = 1$.

\textbf{Problem 5 (Flawed):} Linearization of $f(x,y) = x/y$ at $(2, 1)$.
\textit{Flawed Solution:} $f_x = 1/y = 1$. $f_y = 1/x = 0.5$.
$L = 2 + 1(x-2) + 0.5(y-1)$.
\textbf{Error:} $f_y$ derivative is $-x/y^2$ (quotient/power rule), not $1/x$.
\textbf{Correction:} $f_y(2,1) = -2/1^2 = -2$. Term should be $-2(y-1)$.

\end{document}